% Focus: Package arithmetic torsion and the ACS triple identity with precise hypotheses and a complete proof.

\subsection{Arithmetic torsion as commutator defect}\label{subsec:acs-torsion-def}
% Focus: Define torsion in the discrete path language; clarify terminology vs. classical torsion/curvature.

Arithmetic torsion is the first invariant measuring that two evaluation histories with the same total additive and multiplicative charges need not agree as arrows in the path groupoid of Section~\ref{subsec:flow-threadlike}. In this paper, it is defined as a commutator defect rather than as torsion of an affine connection.

Let
\[
\gamma=(g_1,\dots,g_n)
\]
be a finite arithmetic path, applied in that order, where each step $g_j$ is either an additive update $g_j=\oplus_{\mu_j}:x\mapsto x+\mu_j$ or a multiplicative update $g_j=\otimes_{\lambda_j}:x\mapsto e^{\lambda_j}x$. The evaluation of $\gamma$ at an initial value $x$ is
\[
\nu_x(\gamma)\coloneqq g_n\circ\cdots\circ g_1(x).
\]
We write $\bar\gamma=(g_n,\dots,g_1)$ for the reversed history.

\begin{definition}[Global arithmetic torsion]\label{def:acs-global-torsion}
The \emph{global arithmetic torsion} of $\gamma$ is
\begin{equation}\label{eq:acs-tau-alg}
\tau_{\mathrm{alg}}(\gamma)\coloneqq \nu_x(\gamma)-\nu_x(\bar\gamma).
\end{equation}
\end{definition}

At first sight \eqref{eq:acs-tau-alg} depends on the initial value $x$. The point of the ACS is that this dependence cancels: torsion depends only on the history, not on the basepoint.

\subsection{The Accumulative Commutative Space (ACS)}\label{subsec:acs-definition}
% Focus: Define the ACS coordinates and measures/weights; specify regularity assumptions on paths/domains.

The accumulative commutative space is the plane of total charges. It forgets the current assignment value and remembers only the accumulated additive and multiplicative charges along a history.

\begin{definition}[ACS path]\label{def:acs-path}
The \emph{ACS} is the plane with coordinates $(A,M)$, where $A$ records accumulated additive charge and $M$ records accumulated logarithmic multiplicative charge. For each step $g_j$ of $\gamma$ define
\[
(\Delta A_j,\Delta M_j)=
\begin{cases}
(\mu_j,0), & g_j=\oplus_{\mu_j},\\[4pt]
(0,\lambda_j), & g_j=\otimes_{\lambda_j}.
\end{cases}
\]
Starting from $(A_0,M_0)=(0,0)$, set
\[
(A_k,M_k)\coloneqq\sum_{j=1}^k(\Delta A_j,\Delta M_j),
\qquad 1\le k\le n.
\]
The broken line joining $(A_0,M_0),(A_1,M_1),\dots,(A_n,M_n)$ is the \emph{ACS path} of $\gamma$ and is denoted by $C_\gamma$.
\end{definition}

The endpoint $(A_n,M_n)$ depends only on the total charges and not on their order. Hence $C_\gamma$ and $C_{\bar\gamma}$ always have the same initial and terminal points. When all charges are nonnegative, these two monotone paths bound a canonical region between them. In general we use an oriented filling.

\begin{definition}[ACS filling]\label{def:acs-filling}
An \emph{ACS filling} for $\gamma$ is an oriented piecewise smooth $2$-chain $\Sigma_\gamma$ in the $(A,M)$-plane such that
\begin{equation}\label{eq:acs-filling-boundary}
\partial\Sigma_\gamma=C_{\bar\gamma}-C_\gamma
\end{equation}
as oriented $1$-chains.
\end{definition}

All integrals in this section are understood in the standard piecewise smooth sense. Changing the parameterization of the edges of $C_\gamma$ does not change the line integrals, and changing the parameterization of $\Sigma_\gamma$ does not change the area integrals.

\subsection{Weighted area and boundary formulations}\label{subsec:acs-area-boundary}
% Focus: Derive the weighted area integral and boundary integral expressions with full orientation conventions.

The ACS does not evaluate an expression by itself; it records exactly the commutative data needed to compare two different histories of the same total charge. The basic observation is that evaluation is a weighted boundary integral on the reversed ACS path.

\begin{proposition}[Evaluation as a weighted edge integral]\label{prop:acs-evaluation-formula}
For every arithmetic path $\gamma$ and every initial value $x$,
\begin{equation}\label{eq:acs-evaluation-formula}
\nu_x(\gamma)=e^{M_\gamma}x+\int_{C_{\bar\gamma}} e^M\,dA,
\end{equation}
where $M_\gamma$ is the total multiplicative charge of $\gamma$. Consequently, $\tau_{\mathrm{alg}}(\gamma)$ is independent of $x$.
\end{proposition}

\begin{proof}
Let $I_A\subseteq\{1,\dots,n\}$ be the set of additive indices, and for each $j\in I_A$ define the \emph{future multiplicative charge}
\[
M_j^+\coloneqq\sum_{\substack{k>j\\ g_k=\otimes_{\lambda_k}}}\lambda_k.
\]
When $\gamma$ is evaluated, the additive contribution $\mu_j$ is subsequently scaled by every later multiplicative step, hence by the factor $e^{M_j^+}$. Therefore
\[
\nu_x(\gamma)=e^{M_\gamma}x+\sum_{j\in I_A}\mu_j e^{M_j^+}.
\]
Now traverse the reversed history $\bar\gamma$. The edge corresponding to the original step $g_j=\oplus_{\mu_j}$ is traversed exactly when the current ACS height is $M_j^+$, because all multiplicative steps that were later in $\gamma$ have become earlier in $\bar\gamma$. Hence
\[
\int_{C_{\bar\gamma}} e^M\,dA=\sum_{j\in I_A}\mu_j e^{M_j^+},
\]
with the sign of $\mu_j$ incorporated by the orientation of the edge integral. This gives \eqref{eq:acs-evaluation-formula}.

Applying the same argument to $\bar\gamma$ yields
\[
\nu_x(\bar\gamma)=e^{M_\gamma}x+\int_{C_{\gamma}} e^M\,dA,
\]
since $M_{\bar\gamma}=M_\gamma$. Subtracting cancels the $x$-term, so $\tau_{\mathrm{alg}}(\gamma)$ is independent of $x$.
\end{proof}

Two features are worth recording. First, the kernel $e^M$ is forced: every additive step feels all multiplicative steps that come after it. Second, the boundary expression makes torsion a purely commutative quantity on the ACS.

Define the canonical $1$-form on the ACS by
\begin{equation}\label{eq:acs-eta}
\eta\coloneqq e^M\,dA.
\end{equation}
Then Proposition~\ref{prop:acs-evaluation-formula} implies
\begin{equation}\label{eq:acs-tau-boundary}
\tau_{\mathrm{alg}}(\gamma)=\int_{C_{\bar\gamma}}\eta-\int_{C_\gamma}\eta=\oint_{\partial\Sigma_\gamma}\eta,
\end{equation}
for any ACS filling $\Sigma_\gamma$.

\subsection{The ACS triple identity}\label{subsec:acs-triple-identity}
% Focus: State the triple identity as a theorem; give a complete proof; list invariance properties under reparametrization.

\begin{theorem}[ACS triple identity]\label{thm:acs-triple-identity}
Let $\eta=e^M\,dA$ be the canonical $1$-form on the ACS. For every arithmetic path $\gamma$ and every ACS filling $\Sigma_\gamma$,
\begin{equation}\label{eq:acs-triple-identity}
\tau_{\mathrm{alg}}(\gamma)
=\oint_{\partial\Sigma_\gamma}\eta
=\iint_{\Sigma_\gamma} d\eta
=\iint_{\Sigma_\gamma} e^M\,dM\wedge dA.
\end{equation}
In particular, when all charges are nonnegative and $\Sigma_\gamma$ is the canonical region between the two monotone ACS paths, $\tau_{\mathrm{alg}}(\gamma)$ is the positive $e^M$-weighted area between $C_\gamma$ and $C_{\bar\gamma}$.
\end{theorem}

\begin{proof}
Equation \eqref{eq:acs-tau-boundary} gives the boundary integral expression. Since $d\eta=e^M\,dM\wedge dA$, Stokes' theorem on the plane yields
\[
\oint_{\partial\Sigma_\gamma}\eta=\iint_{\Sigma_\gamma} d\eta
=\iint_{\Sigma_\gamma} e^M\,dM\wedge dA,
\]
proving \eqref{eq:acs-triple-identity}.
\end{proof}

The right-hand side depends only on the oriented $2$-chain $\Sigma_\gamma$ with boundary \eqref{eq:acs-filling-boundary}; in particular it is invariant under reparameterizations of the boundary paths and of the filling.

From an algebraic viewpoint, the ACS is the commutative shadow of the noncommutative path calculus: the charge map $\gamma\mapsto(A_n,M_n)$ is an abelianization of histories, while the torsion is the defect detected by comparing a word to its reverse.
